\setlength{\baselineskip}{20pt}
\chapter{面向无额外推理成本的特征增强算法研究（DenoiseRep）}
\label{cha:chap3}

\section{研究问题与整体框架设计}
本章提出基础方法 DenoiseRep，目标是在不增加推理开销的前提下提升判别特征质量。核心思路是将骨干网络的级联嵌入层解释为逐层去噪过程，使特征提取与去噪增强在同一主干内协同完成。与仅在最终输出端附加去噪器不同，该设计能够在层级语义空间中逐步抑制噪声累积。

\section{联合特征提取与特征去噪}

\subsection{扩散过程建模}
设骨干网络共有 $N$ 层嵌入映射，第 $n$ 层特征记为 $\mathbf{h}_n$。本文将各层特征视作扩散建模的样本，构造前向加噪过程，并以层级与时间步对齐的方式设定噪声调度，从而在特征提取流程中引入逐层去噪视角。

\subsection{去噪网络训练}
本文在每层后引入去噪分支 $G_n$，训练阶段冻结主干参数，仅优化去噪分支。其噪声预测损失写为
\begin{equation}
\mathcal{L}_{denoise}=\mathbb{E}_{\mathbf{h},\epsilon,t}\left[\left\|\epsilon-G_n(\mathbf{h}_{n,t},t)\right\|_2^2\right].
\end{equation}
其中 $\mathbf{h}_{n,t}$ 为第 $n$ 层特征在扩散步 $t$ 下的噪声样本，$\epsilon$ 为高斯噪声。该训练方式同时兼容两类范式：
\begin{enumerate}
  \item \textbf{label-free}：仅优化去噪损失，利用无标签数据学习特征噪声分布；
  \item \textbf{label-augmented}：将任务损失与去噪损失加权联合，兼顾判别监督与噪声抑制。
\end{enumerate}

\section{特征提取与去噪的参数融合}

\subsection{双分支结构设计}
训练阶段在每层引入去噪分支，与原嵌入层形成双分支结构；推理阶段通过参数融合恢复单分支，实现训练增强与部署高效的统一。

\subsection{参数融合数学推导}
为消除推理阶段额外开销，本文将去噪分支参数折叠进原嵌入层参数。设原层映射为
\begin{equation}
\mathbf{y}=\mathbf{W}\mathbf{x}+\mathbf{b},
\end{equation}
融合后得到等价映射
\begin{equation}
\mathbf{y}'=\mathbf{W}'\mathbf{x}+\mathbf{b}'.
\end{equation}
其中 $\mathbf{W}'$、$\mathbf{b}'$ 吸收了去噪分支线性修正项。融合后网络保持单分支结构，参数量与计算图复杂度不增加，因而可实现 computation-free 的部署增强。

\section{本章实验设计}
\subsection{实验设置与评价指标}
本章实验聚焦基础方法有效性，采用 Market-1501、DukeMTMC-reID、MSMT17、CUHK-03 四个 ReID 数据集，评价指标为 mAP 与 Rank-1。

\subsection{实现细节与训练策略}
骨干采用 TransReID-SSL 配置，并在相同训练预算下与基线比较。\,[需要补充具体训练超参与实现细节]

\section{实验结果与分析}
\subsection{基础效果与对比分析}
表\ref{tab:base-reid}显示，DenoiseRep 在四个 ReID 数据集上均实现稳定增益，说明在不改变骨干结构的前提下，特征层去噪能够有效提升表征判别性。

\begin{table}[!htb]
\centering
\caption{基础方法在 ReID 数据集上的结果}
\label{tab:base-reid}
\small
\begin{tabular}{|c|c|c|c|}
\hline
数据集 & 方法 & mAP(\%) & Rank-1(\%) \\
\hline
DukeMTMC-reID & Baseline & 80.40 & 87.80 \\
DukeMTMC-reID & +DenoiseRep & 80.76 & 88.22 \\
\hline
MSMT17 & Baseline & 66.30 & 84.80 \\
MSMT17 & +DenoiseRep & 66.81 & 85.34 \\
\hline
Market-1501 & Baseline & 91.20 & 95.80 \\
Market-1501 & +DenoiseRep & 91.59 & 96.21 \\
\hline
CUHK-03 & Baseline & 83.50 & 85.90 \\
CUHK-03 & +DenoiseRep & 83.69 & 86.02 \\
\hline
\end{tabular}
\end{table}

进一步实验表明，在引入标签监督后，方法相较仅监督基线仍有额外提升，体现了“监督判别信号+无监督去噪信号”的互补性。

在 CNN 与 Transformer 两类主流骨干上，基础方法均呈现一致增益。尤其在中小容量骨干中，提升幅度更明显，说明特征噪声抑制对弱表征模型具有更高边际收益。

\subsection{效率与开销验证}
推理效率实验显示：仅在末层附加去噪分支会引入额外时延，而逐层去噪并执行参数融合后，推理时间与基线基本一致，同时精度保持提升。该结果从工程层面验证了融合机制的有效性。

\section{本章小结}
本章提出了基础方法 DenoiseRep，并在本章内完成了有效性、效率与对比实验验证。结果表明，分层去噪与参数融合可在不增加推理开销的前提下稳定提升判别特征质量。
